\section{Grouped Learning and Integrity Publication}
\label{app:grouped-learning}

\subsection{Retrospective SIM Training Contract}
The SIM11--19 registry contains 176 observation records from 88 physical
cases. Connected components join shared physical histories and observation
replicas before splitting. Of 56 dependency groups, \MLTrainGroups{} are
used for fitting, \MLValGroups{} for model selection and \MLHoldGroups{}
for the reserved development evaluation: \MLTrainRecords{},
\MLValRecords{} and \MLHoldRecords{} records, respectively. This implements
the requested approximately 70/30 development/reserved partition with an
internal validation subset. All campaigns were already open; the split is
retrospective and does not establish unseen-source or external generalization.

Both ExtraTrees variants use 400 trees, minimum leaf size two, square-root
feature sampling, balanced class weights and seed 2601026. Fitting uses
every third report starting at index 360. The target is normal outside an
intervention and one of 18 load/line sources during it; recovery is not a
separate training class. The baseline has 846 trailing features; SHORT adds
611 features. Their input is the same 47-channel representation with a
300-report empirical reference, colored filtering and median/MAD scaling.
No DAE fit or RAW annotation enters fitting. The empirical reference has no
startup stationarity guarantee; it is distinct from the physical filter's
admission gate.\AIDraftingCredit{}

The publication adapter requires six consecutive winning source labels.
Scores are not claimed to be calibrated probabilities. Selection maximizes
validation diagnoses with correct family and source emitted before OFF,
then minimizes extra alarms, preferring the baseline on a complete tie.
SHORT is frozen after \MLShortValTimely{}/\MLShortValEvents{} versus
\MLBaseValTimely{}/\MLBaseValEvents{} validation successes. There is no
refit on validation and no subsequent parameter or threshold selection
from the reserved \MLHoldRecords{} records. The declared metric associates
the first temporal alarm independently of its predicted class, within
30 seconds or before the next onset, whichever comes first; success also
requires diagnosis before the current OFF.

The reserved result is \MLHoldTimely{}/\MLHoldEvents{}, comprising
\MLHoldLoad{}/\MLHoldLoadTotal{} load and
\MLHoldLine{}/\MLHoldLineTotal{} line interventions, with \MLHoldExtras{}
extra alarms. Only two distinct normal dependency groups exist: one in
training and one reserved; validation has none. The
\MLHoldNormalSeconds{} normal seconds are repeated observations of the
reserved normal group, with \MLHoldNormalAlarms{} alarms. They do not
provide independent long-duration operational exposure. In particular,
neither report-level binomial intervals nor a validation false-alarm rate
would be justified. Load weakness and the common F01 ancestry limit
transfer claims even when the group assignment is disjoint.

The selected model is also replayed on all twelve H04 records, using the
same successor event metrics as the physical controls. It obtains
\ETCommonTimely{}/\ETCommonTotal{} timely decisions, \ETCommonExtras{}
extras, \ETCommonOrphan{} orphan restorations and \ETCommonMissing{}
missing closures, including missed events. All missed events are loads;
all tested lines are diagnosed before OFF. Its six-argmax publisher has no
confidence threshold, so this is not an identical decision-clock ablation.
The three H04 event groups are disjoint from the SIM fitting, validation
and reserved groups, but its single normal group is shared with the
reserved group. The inference-only replay reads no truth, registry or split.
An earlier roster contained unused event/normal metadata; it is preserved
and documented. Removing that metadata reproduces every output array and
event exactly on all twelve records, before separate post-seal scoring.

\subsection{Missing-PMU Notice Timing}
The continuous RAW replay contains \MissingEpisodes{} observed missing
intervals. The exact submitted artifact and the immediate availability
interface both identify all intervals and their affected PMU. First correct
notices have mean delay \MissingLegacyDelay{} versus
\MissingImmediateDelay{} seconds, respectively. The clock is the received
report timestamp relative to the first affected report; zero means emission
on that report, not zero computation or communication latency. Notices
before onset cannot match a later missing interval, and a later correct
source cannot repair an incorrect first notice. This is a publication benefit
from already-observed flags and finite-value masks, not evidence of a new
malicious-attack detector or improved physical-source inference.

\subsection{Segmented ML Transfer and Source Composition}
The classifier consumes the phase-A voltage/current magnitudes and angles,
frequency and ROCOF in the declared eight-PMU order. Complete segments
require all 48 measurements finite and all presence flags true; missing
reports are never removed to join their neighboring samples. Each segment
gets its own 300-report empirical reference and remains unavailable through
index 359. Segment startup is not certified normal and physical state is
not preserved across gaps. The same frozen classifier and publication rule
are used; no RAW source label, event code or score enters inference.

For a submitted-artifact prediction of line or load, including its predicted
physical component in a mixed missing case, source replacement requires the
first ML diagnosis of an episode with that same predicted family, emitted
inside the completed window no later than the artifact's own emission.
Later revisions cannot repair that first diagnosis. The earliest eligible
diagnosis is selected; otherwise the ZIP source is retained. Fault,
generation, integrity outputs and all window class labels and times are
unchanged. The archived Top-3 list remains identified as the original ZIP
list; it is not presented as the hybrid's reranked posterior.

The completed replay has \RawHybridSegments{} complete segments and
\RawHybridRoutes{} selected source substitution, with
\RawHybridFallbacks{} fallbacks. The evaluator verifies full report
coverage, explicit startup unavailability, causal message times and every
replacement against the frozen rule before scoring. Source identity is
supplied only to this later evaluator by the original metadata workbook:
fault BUS39, line 23--24, generation BUS2 and load BUS7. Current event-code
support determines which fixed windows contain each family; the workbook's
approximate times do not supply exact action onsets or restorations.
All physical-positive windows, including generation outside the learned
catalog, enter the denominator. Paired joint correctness is
\RawZIPJoint{}/\RawZIPTotal{} versus
\RawHybridJoint{}/\RawHybridTotal{}. These are dependent windows of one
retrospective record; no population-level significance or official
competition improvement is inferred.
